\noindent
{\Large \textbf{\section{Особенности реализации}}} \par

Программа реализована на языке C (файл cache.c) и компилируется с оптимизацией \texttt{-O2}. Все замеры основаны на обходе цепочки ячеек в массиве: в каждой ячейке хранится смещение следующей, а цепочка замкнута в кольцо. Параметры замера задаются константами: \texttt{ACCESSES\_PER\_RUN} - число обращений в одном замере, \texttt{RUN\_COUNT} - число повторных замеров, по которым вычисляется матожидание. Переменная \texttt{result} (тип \texttt{volatile}) используется только для того, чтобы компилятор не удалил цикл обхода. Ниже описаны функции программы.

\vspace{8pt}

\textbf{Функция time\_seconds } \par
\textbf{Вход:} параметров нет, используется текущее показание системных монотонных часов \texttt{CLOCK\_MONOTONIC} \par
\textbf{Выход:} double - текущее время в секундах (число секунд с некоторого момента отсчета, имеет смысл только разность двух значений) \par
\textbf{Описание:} Возвращает показание монотонных часов \texttt{CLOCK\_MONOTONIC}, приведенное к секундам. Монотонные часы не зависят от изменения системного времени, поэтому разность двух показаний корректно задает длительность замера.\par
\vspace{10pt}

\begin{lstlisting}[caption=time\_seconds, label=lst:time-seconds]
static double time_seconds(void) {
    struct timespec ts;
    clock_gettime(CLOCK_MONOTONIC, &ts);
    return ts.tv_sec + ts.tv_nsec * 1e-9;
}
\end{lstlisting}

\textbf{Функция expected\_value } \par
\textbf{Вход:} const double *values - массив результатов замеров; int count - число замеров \par
\textbf{Выход:} double - матожидание (среднее арифметическое) значений массива \par
\textbf{Описание:} Каждый замер считается равновероятным (вероятность исхода равна $1 / count$), поэтому матожидание вычисляется как сумма значений, деленная на их число. Значение используется как итоговое время обращения для одной точки измерения.\par
\vspace{10pt}

\begin{lstlisting}[caption=expected\_value, label=lst:expected-value]
static double expected_value(const double *values, int count) {
    double sum = 0;
    for (int i = 0; i < count; i++)
        sum += values[i];
    return sum / count;
}
\end{lstlisting}

\textbf{Функция build\_ring\_by\_memory\_step } \par
\textbf{Вход:} char *buffer - массив; size\_t buffer\_size - размер рабочей части массива в байтах; size\_t memory\_step\_bytes - шаг между ячейками цепочки в байтах \par
\textbf{Выход:} изменяется содержимое массива \texttt{buffer}: в нем построена замкнутая цепочка ячеек \par
\textbf{Описание:} В ячейку по смещению $i$ записывается смещение следующей ячейки $(i + memory\_step\_bytes) \bmod buffer\_size$. Заполняются только ячейки, расположенные через шаг друг от друга. Операция взятия остатка замыкает цепочку в кольцо, что позволяет обходить массив неограниченное число раз.\par

\newpage

\begin{lstlisting}[caption=build\_ring\_by\_memory\_step, label=lst:build-ring-by-memory-step]
static void build_ring_by_memory_step(char *buffer, size_t buffer_size, size_t memory_step_bytes) {
    // заполняем только ячейки, которые будем посещать (по одной на каждые memory_step_bytes байт)
    for (size_t i = 0; i < buffer_size; i += memory_step_bytes)
        *(size_t *)(buffer + i) = (i + memory_step_bytes) % buffer_size;
}
\end{lstlisting}

\textbf{Функция build\_ring\_chaotic } \par
\textbf{Вход:} char *buffer - массив; size\_t buffer\_size - размер рабочей части массива в байтах (степень двойки); size\_t line\_bytes - размер строки кэша в байтах \par
\textbf{Выход:} изменяется содержимое массива \texttt{buffer}: в нем построена замкнутая цепочка, посещающая все строки в хаотичном порядке \par
\textbf{Описание:} Массив делится на строки по $line\_bytes$ байт, их число $m$ равно \\ $buffer\_size / line\_bytes$. В начале строки с номером $i$ записывается смещение начала строки с номером $(5i + 1) \bmod m$. Так как $m$ является степенью двойки, а числа 5 и 1 удовлетворяют условиям полного периода линейного конгруэнтного генератора, цепочка проходит по каждой строке ровно один раз. Порядок обхода не имеет постоянного шага, поэтому аппаратный префетчер не может предсказать следующий адрес.\par
\vspace{10pt}

\begin{lstlisting}[caption=build\_ring\_chaotic, label=lst:build-ring-chaotic]
static void build_ring_chaotic(char *buffer, size_t buffer_size, size_t line_bytes) {
    size_t line_count = buffer_size / line_bytes;
    // в начале каждой строки храним смещение начала следующей строки
    for (size_t i = 0; i < line_count; i++)
        *(size_t *)(buffer + i * line_bytes) = ((5 * i + 1) % line_count) * line_bytes;
}
\end{lstlisting}

\textbf{Функция walk } \par
\textbf{Вход:} char *buffer - массив; size\_t buffer\_size - размер рабочей части в байтах; size\_t memory\_step\_bytes - шаг цепочки (или размер строки для хаотичного порядка); \\ int isHaosStep - режим построения цепочки: 1 - хаотичный порядок, 0 - постоянный шаг \par
\textbf{Выход:} double - среднее время одного обращения в наносекундах (матожидание по \texttt{RUN\_COUNT} замерам) \par
\textbf{Описание:} Функция строит цепочку выбранным способом и выполняет прогрев: пять полных проходов по цепочке, чтобы данные закрепились в кэше. Затем \texttt{RUN\_COUNT} раз выполняется замер: фиксируется время, выполняется \texttt{ACCESSES\_PER\_RUN} обращений вида \texttt{offset = *(size\_t *)(buffer + offset)}, и разность времен делится на число обращений и переводится в наносекунды. Адрес каждого обращения берется из ячейки, прочитанной предыдущим обращением, поэтому процессор не может выполнять обращения параллельно, и измеряется именно задержка того уровня памяти, где находятся данные. Запись конечного смещения в переменную \texttt{result} объявленную как \texttt{volatile} не позволяет компилятору удалить цикл как не имеющий результата. Итоговое значение вычисляется функцией \texttt{expected\_value} как матожидание результатов замеров.\par
\vspace{10pt}

\newpage
\begin{lstlisting}[caption=walk, label=lst:walk]
static double walk(char *buffer, size_t buffer_size, size_t memory_step_bytes, int isHaosStep) {
    double run_times_ns[RUN_COUNT];
    if (isHaosStep) build_ring_chaotic(buffer, buffer_size, memory_step_bytes);
    else build_ring_by_memory_step(buffer, buffer_size, memory_step_bytes);
    size_t offset = 0;
    // прогрев
    for (size_t i = 0; i < 5 * (buffer_size / memory_step_bytes) + 1; i++)
        offset = *(size_t *)(buffer + offset);
    for (int run = 0; run < RUN_COUNT; run++) {
        double start_time = time_seconds();
        // следующий адрес читаем из только что прочитанной ячейки
        for (long access = 0; access < ACCESSES_PER_RUN; access++)
            offset = *(size_t *)(buffer + offset);
        run_times_ns[run] = (time_seconds() - start_time) / ACCESSES_PER_RUN * 1e9;
    }
    result = offset;
    return expected_value(run_times_ns, RUN_COUNT);
}
\end{lstlisting}

\textbf{Функция main } \par
\textbf{Вход:}  размер буфера 256 МБ, размер массива в тесте строки 128 МБ, шаг 256 КБ и число ячеек до 32 в тесте ассоциативности, константы \texttt{ACCESSES\_PER\_RUN} и \texttt{RUN\_COUNT} \par
\textbf{Выход:} int - код завершения программы \par
\textbf{Описание:} Выделяется массив максимального размера 256 МБ. Далее последовательно выполняются три эксперимента. В первом шаг цепочки увеличивается от 8 до 512 байт при размере массива 128 МБ; размером строки считается тот шаг, при переходе к которому время обращения возрастает относительно предыдущего шага сильнее всего. Во втором размер массива удваивается от 4 КБ до 256 МБ, цепочка строится в хаотичном порядке с найденным размером строки. В третьем берется $k$ ячеек от 1 до 32 с шагом 256 КБ, которые попадают в один набор кэша.\par
\vspace{10pt}

\newpage
\begin{lstlisting}[caption=main, label=lst:main]
int main(void) {
    // зададин максимальный размер массива в 256МБ
    size_t max_buffer_size = 256UL << 20;
    char *buffer = malloc(max_buffer_size);

    // 1. Размер строки
    printf("1. Размер строки\n");
    printf("шаг, байт   нс/обращение\n");
    size_t line_bytes = 0;
    double prev_time = 0, max_growth = 0;
    for (size_t memory_step = 8; memory_step <= 512; memory_step *= 2) {
        double time = walk(buffer, 128UL << 20, memory_step, 0);
        printf("%9zu   %8.3f\n", memory_step, time);
        // строка там, где время обращения вырастает сильнее всего
        if (prev_time > 0 && time / prev_time > max_growth) {
            max_growth = time / prev_time;
            line_bytes = memory_step;
        }
        prev_time = time;
    }
    printf("Размер строки: %zu байт\n", line_bytes);

    // 2. Размер кэша. 
    // Размер строки берем из предыдущего замера.
    printf("\n2. Размер кэша\n");
    printf("размер, КБ   нс/обращение\n");
    for (size_t size = 4096; size <= max_buffer_size; size *= 2)
        printf("%10zu   %8.3f\n", size / 1024, walk(buffer, size, line_bytes, 1));

    // 3. Ассоциативность. 
    printf("\n3. Ассоциативность\n");
    printf("  k   нс/обращение\n");
    for (size_t elements = 1; elements <= 32; elements++)
        printf("%3zu   %8.3f\n", elements, walk(buffer, elements * (256UL << 10), 256UL << 10, 0));

    free(buffer);
    return 0;
}
\end{lstlisting}

\newpage
